\documentclass[12pt,a4paper]{article}

% ------------------------------------------------
% PACKAGES
% ------------------------------------------------
\usepackage[margin=1in]{geometry}
\usepackage{setspace}
\usepackage{titlesec}
\usepackage{fancyhdr}
\usepackage{hyperref}
\usepackage{booktabs}
\usepackage{longtable}
\usepackage{lastpage}
\usepackage{enumitem}
\usepackage{array}

\setstretch{1.15}

% ------------------------------------------------
% HEADER / FOOTER
% ------------------------------------------------
\pagestyle{fancy}
\fancyhf{}
\lhead{IBDP Psychology (First Assessment 2025)}
\rhead{Paper 2 Framework}
\rfoot{Page \thepage\ of \pageref{LastPage}}

\hypersetup{
    colorlinks=true,
    linkcolor=black
}

% ------------------------------------------------
% DOCUMENT
% ------------------------------------------------
\begin{document}

\begin{center}
    {\Large \textbf{IBDP Psychology}}\\[6pt]
    {\large Paper 2 – Structure, Instructions and Assessment Framework}\\[10pt]
    \rule{\textwidth}{0.4pt}
\end{center}

\vspace{0.5cm}

\section*{Purpose of This Document}

This document provides an internally developed summary of the structure and expectations of Paper 2 under the First Assessment 2025 specification. It is intended for academic planning and instructional clarity. This guide paraphrases structural requirements and does not reproduce proprietary examination content.

\newpage

% ------------------------------------------------
\section{Paper 2 Overview}

\begin{itemize}
    \item \textbf{Duration:} 1 hour 30 minutes
    \item \textbf{Total Marks:} 30
    \item \textbf{Weighting:}
    \begin{itemize}
        \item Standard Level (SL): 25\%
        \item Higher Level (HL): 25\%
    \end{itemize}
\end{itemize}

\subsection*{Assessment Focus}

Paper 2 evaluates:

\begin{itemize}
    \item In-depth understanding of one optional theme
    \item Ability to construct a sustained analytical essay
    \item Integration of research evidence
    \item Critical evaluation of psychological explanations
    \item Logical and coherent argument development
\end{itemize}

\subsection*{Content Scope}

Paper 2 assesses content from the optional themes studied by the school. Examples of optional themes include:

\begin{itemize}
    \item Abnormal psychology
    \item Developmental psychology
    \item Health psychology
    \item Psychology of human relationships
\end{itemize}

Students are assessed only on optional content that has been taught.

\newpage

% ------------------------------------------------
\section{Paper 2 Structure}

Paper 2 consists of two sections: Section A and Section B.

\subsection*{Total Requirement}

\begin{itemize}
    \item Two questions are provided.
    \item Students answer one question only.
    \item Each question is worth 30 marks.
\end{itemize}

Both questions require extended analytical responses.

\newpage

% ------------------------------------------------
\section{Section A}

\subsection*{Question Format}

\begin{itemize}
    \item Extended-response essay
    \item Focused within one optional theme
    \item Uses higher-order command terms
\end{itemize}

\subsection*{Typical Cognitive Demand}

Questions may require students to:

\begin{itemize}
    \item Discuss psychological explanations
    \item Evaluate one or more theories or models
    \item Examine research evidence
    \item Assess the strengths and limitations of approaches
\end{itemize}

\subsection*{Skills Assessed}

\begin{itemize}
    \item Knowledge and understanding (AO1)
    \item Synthesis and evaluation (AO3)
    \item Integration of research
\end{itemize}

\subsection*{High-Quality Response Characteristics}

A strong Section A response should:

\begin{itemize}
    \item Clearly define relevant terminology
    \item Explain theories accurately and coherently
    \item Incorporate relevant research appropriately
    \item Demonstrate balanced evaluation
    \item Present a sustained and logical argument
    \item Conclude with a reasoned judgement
\end{itemize}

\newpage

% ------------------------------------------------
\section{Section B}

\subsection*{Question Format}

\begin{itemize}
    \item Extended-response essay
    \item Focused within one optional theme
    \item Framed using evaluative or analytical command terms
\end{itemize}

Section B offers an alternative essay prompt to Section A. Students select one question from either Section A or Section B.

\subsection*{Typical Cognitive Demand}

Questions in Section B may require:

\begin{itemize}
    \item Comparative analysis
    \item Critical evaluation of research methodology
    \item Examination of competing explanations
    \item Application of theory to broader implications
\end{itemize}

\subsection*{Skills Assessed}

\begin{itemize}
    \item Conceptual understanding
    \item Analytical reasoning
    \item Evaluation of evidence
    \item Structured argumentation
\end{itemize}

\newpage

% ------------------------------------------------
\section{Mark Allocation and Assessment Emphasis}

\begin{longtable}{p{3cm} p{3cm} p{7cm}}
\toprule
\textbf{Section} & \textbf{Marks} & \textbf{Assessment Emphasis} \\
\midrule
Section A & 30 & Extended analytical essay requiring explanation and evaluation \\
Section B & 30 & Alternative extended analytical essay requiring critical reasoning \\
\bottomrule
\end{longtable}

Only one question is answered. Maximum mark available: 30.

\newpage

% ------------------------------------------------
\section{Recommended Essay Structure}

A structured response may include:

\begin{enumerate}[leftmargin=1.5cm]
    \item Introduction
    \begin{itemize}
        \item Define key terms
        \item Clarify focus
        \item Outline argument
    \end{itemize}
    \item Theoretical Explanation
    \item Research Integration
    \item Critical Analysis
    \begin{itemize}
        \item Strengths
        \item Limitations
        \item Methodological considerations
        \item Cultural or ethical implications (where relevant)
    \end{itemize}
    \item Conclusion
    \begin{itemize}
        \item Directly address command term
        \item Provide reasoned judgement
    \end{itemize}
\end{enumerate}

\newpage

% ------------------------------------------------
\section{Pedagogical Implications}

\subsection*{Content Preparation}

\begin{itemize}
    \item Deep mastery of at least one optional theme
    \item Clear conceptual mapping of theories
\end{itemize}

\subsection*{Essay Training}

\begin{itemize}
    \item Command term analysis
    \item Timed 90-minute simulations
    \item Evaluation paragraph drills
\end{itemize}

\subsection*{Research Integration}

\begin{itemize}
    \item Analytical use of studies (not narrative description)
    \item Linking evidence directly to argument
\end{itemize}

\vfill

\begin{center}
\rule{\textwidth}{0.4pt} \\
\small{This document is an internally developed academic resource summarizing structural expectations of IB Psychology Paper 2 (First Assessment 2025). It is not affiliated with or endorsed by the International Baccalaureate Organization.}
\end{center}

\end{document}